%% kapa: exam=1 round=20 year=2009 period=1 form=A subject=economics subject_ko=경제원론 seq=29 number=69
%% source: Q-Net 감정평가사 기출문제 (cst00309) artlSeq=5024178 file=extracted/1008/5024178/2009년 제20회 감정평가사 1차시험 기출문제/1교시.hwp
%% conversion: kordoc 4.24.51 -> tools/gen1.py (auto); review=proofread-sonnet; flags=none
\begin{question}{69}
생산함수 $Q= { 3K+2L }$ 이 나타내는 생산기술의 특성으로 옳은 것은? (단, $Q$ 는 생산량, $K$ 는 자본의 투입량, $L$ 은 노동의 투입량이다.)
\begin{choices}
  \item 요소 간 완전대체
  \item 자본의 한계생산 체증
  \item 노동의 한계생산 체증
  \item 한계기술대체율 체감
  \item 규모에 대한 수익체증
\end{choices}
\end{question}
\answer{1}
